% ============================================================================
\section{Giới thiệu}
\label{sec:intro}
% ============================================================================

\subsection{Đặt vấn đề}

\textbf{Đọc hiểu máy trích xuất} (extractive Machine Reading Comprehension --- MRC) là bài
toán: cho một đoạn văn và một câu hỏi, chỉ ra đoạn chữ \emph{trong chính văn bản đó} trả
lời câu hỏi, hoặc kết luận rằng văn bản không chứa câu trả lời. Đây là thành phần ``bộ
đọc'' của các hệ thống hỏi đáp trên tài liệu và của các pipeline truy hồi--sinh hiện nay.
Với tiếng Việt, bộ dữ liệu UIT-ViQuAD~2.0 \cite{nguyen2022vlsp} cung cấp nền tảng đánh giá
chung theo định dạng SQuAD~2.0 \cite{rajpurkar2018know}, gồm cả câu hỏi \emph{không có đáp án}.

Hai họ mô hình Transformer tiền huấn luyện thường được dùng cho bài toán này khác nhau ở
một điểm mà môn học nhấn mạnh từ Buổi~3 --- \textbf{tách từ}:

\begin{itemize}
  \item \textbf{PhoBERT} \cite{nguyen2020phobert} là mô hình \emph{đơn ngữ} tiếng Việt,
        tiền huấn luyện trên văn bản \emph{đã tách từ}: các âm tiết của một từ được nối
        bằng dấu gạch dưới (``Hà\_Nội'', ``thủ\_đô''). Đơn vị đầu vào của nó là \emph{từ}.
  \item \textbf{XLM-R} \cite{conneau2020xlmr} là mô hình \emph{đa ngữ} (100 ngôn ngữ), dùng
        SentencePiece \cite{kudo2018sentencepiece} trên văn bản thô. Đơn vị đầu vào của nó
        là \emph{mảnh từ con} (subword), không cần tách từ trước.
\end{itemize}

Khác biệt này không chỉ là chi tiết cài đặt. Một mô hình đọc hiểu \emph{trích xuất} trả
lời bằng cách chọn token bắt đầu và token kết thúc. Nếu đơn vị nhỏ nhất là một \emph{từ đã
tách}, mô hình không thể trả về nửa từ: khi ranh giới của đáp án vàng cắt ngang một từ ghép
theo bộ tách từ, mô hình word-level \emph{không có cách nào} khớp chính xác. Ngược lại, mô
hình subword có thể mất lợi thế khi đơn vị ngữ nghĩa của tiếng Việt là từ nhiều âm tiết.
Luận điểm xuyên suốt của môn học --- \emph{``phải dựa trên đặc điểm dữ liệu để chọn phương
pháp''} --- vì vậy có một phiên bản cụ thể, kiểm chứng được: \emph{đặc điểm tách từ của
tiếng Việt ảnh hưởng thế nào đến việc chọn giữa mô hình word-level và subword cho MRC?}

Các bảng xếp hạng MRC chỉ báo một điểm tổng EM/F1. Điểm tổng trả lời được câu hỏi ``mô hình
nào cao hơn'' nhưng không trả lời được ``mô hình sai ở đâu, và vì sao''. Đồ án này đặt
trọng tâm vào câu hỏi thứ hai.

\subsection{Phát biểu bài toán}
\label{sec:problem}

Cho đoạn văn $C = (c_1, \dots, c_{|C|})$ là chuỗi ký tự và câu hỏi $Q$. Hệ thống cần tìm
cặp chỉ số ký tự $(s, e)$, $0 \le s < e \le |C|$, sao cho đoạn con $C[s{:}e]$ trả lời $Q$;
hoặc trả về chuỗi rỗng $\varnothing$ nếu $C$ không chứa câu trả lời.

\begin{table}[htbp]
\centering
\caption{Đặc tả đầu vào, đầu ra}
\label{tab:io}
\renewcommand{\arraystretch}{1.2}
\begin{tabularx}{\textwidth}{@{}l L@{}}
\toprule
\textbf{Thành phần} & \textbf{Mô tả} \\
\midrule
Đầu vào & \code{context}: đoạn văn Wikipedia tiếng Việt; \code{question}: câu hỏi tự nhiên \\
Đầu ra & \code{answer}: một chuỗi \emph{con} của \code{context}, hoặc chuỗi rỗng \\
Loại bài toán & Trích xuất (extractive), không sinh văn bản \\
Dữ liệu & UIT-ViQuAD~2.0 (định dạng SQuAD~2.0) \\
Thang đo & Exact Match (EM), F1 trên token; cộng các phép đo chẩn đoán ở Mục~\ref{sec:method} \\
\bottomrule
\end{tabularx}
\end{table}

\subsection{Câu hỏi nghiên cứu}
\label{sec:rq}

Theo hướng dẫn của môn học ở Buổi~3, câu hỏi nghiên cứu được đặt ở dạng kiểm chứng được.

\begin{table}[htbp]
\centering
\caption{Bốn câu hỏi nghiên cứu và phép đo trả lời từng câu}
\label{tab:rq}
\renewcommand{\arraystretch}{1.22}
\small
\begin{tabularx}{\textwidth}{@{}l L L@{}}
\toprule
& \textbf{Câu hỏi} & \textbf{Phép đo (Mục)} \\
\midrule
\textbf{RQ1} & Mô hình word-level (PhoBERT) có nhạy với ranh giới từ hơn mô hình subword (XLM-R) không? & Trần EM do biên từ; điểm trên câu lệch/không lệch biên; loại lỗi biên (\ref{sec:res-rq1}) \\
\textbf{RQ2} & Hiệu năng suy giảm theo độ dài ngữ cảnh và độ dài đáp án có khác nhau giữa hai họ mô hình không? & Tách nhóm theo độ dài (\ref{sec:res-rq2}) \\
\textbf{RQ3} & Mô hình thật sự nhận ra câu không có đáp án, hay chỉ học thiên lệch về việc từ chối? & Mức sàn ``luôn từ chối''; precision/recall của hành vi từ chối (\ref{sec:res-rq3}) \\
\textbf{RQ4} & Transformer có vượt baseline không-neural ở \emph{mọi} nhóm câu hỏi không? & So TF-IDF trên từng nhóm (\ref{sec:res-rq4}) \\
\bottomrule
\end{tabularx}
\end{table}

\textbf{Giả thuyết ban đầu cho RQ1} (ghi trong đề cương trước khi chạy): PhoBERT yếu hơn ở
các câu có ranh giới từ ghép phức tạp vì \emph{lỗi tách từ lan truyền} vào mô hình, còn
XLM-R không phụ thuộc bước này. Hai mô hình được chọn vì chúng khác nhau \emph{đúng ở khâu
tokenization} --- đây là một tương phản được thiết kế, không phải so sánh ngẫu nhiên.
Mục~\ref{sec:res-rq1} đối chiếu giả thuyết này với số liệu.

\subsection{Đối chiếu với yêu cầu môn học}
\label{sec:requirements}

Yêu cầu bắt buộc của đồ án CS221 là giải một bài toán NLP có đầu vào là văn bản hoặc tiếng
nói; sản phẩm là mô hình được huấn luyện (``train model thôi, mọi người không cần làm
app''), kèm báo cáo và thuyết trình. Bảng~\ref{tab:req} ánh xạ các yêu cầu và các chủ đề
được nhấn mạnh trong môn học vào nội dung đồ án.

\begin{table}[htbp]
\centering
\caption{Đối chiếu yêu cầu và nội dung môn học với đồ án}
\label{tab:req}
\renewcommand{\arraystretch}{1.18}
\small
\begin{tabularx}{\textwidth}{@{}L L c@{}}
\toprule
\textbf{Yêu cầu / nội dung môn học} & \textbf{Thực hiện trong đồ án} & \textbf{Mục} \\
\midrule
Bài toán NLP, đầu vào văn bản & Đọc hiểu trích xuất tiếng Việt & \ref{sec:problem} \\
Huấn luyện mô hình & Fine-tune PhoBERT-base-v2 và XLM-R-base trên toàn bộ tập train & \ref{sec:finetune} \\
Buổi 3 --- tiền xử lý, tách từ & Bộ tách từ cho PhoBERT; phép đo tương thích biên từ & \ref{sec:segtok}, \ref{sec:alignment} \\
Buổi 3 --- quy trình nghiên cứu & Câu hỏi nghiên cứu dạng kiểm chứng; giả thuyết ghi trước & \ref{sec:rq} \\
Buổi 4 --- TF-IDF & Baseline truy hồi câu bằng TF-IDF & \ref{sec:baselines} \\
Buổi 5--7 --- phân tích lỗi & Phân loại lỗi tự động trên toàn bộ tập chấm; phân tích định tính & \ref{sec:taxonomy}, \ref{sec:qualitative} \\
Buổi 7 --- giới hạn độ dài, từ vs âm tiết & Cửa sổ trượt; so sánh độ dài chuỗi token hai tokenizer & \ref{sec:windowing}, \ref{sec:res-rq2} \\
Buổi 7 --- chất lượng nhãn & Kiểm định bộ stress-test của nhóm & \ref{sec:audit} \\
\bottomrule
\end{tabularx}
\end{table}

\subsection{Phạm vi}

Đồ án giới hạn trong hỏi đáp trích xuất trên \emph{một} đoạn văn cho trước. Nằm ngoài
phạm vi: hỏi đáp sinh văn bản, truy hồi đoạn văn từ kho lớn, hỏi đáp hội thoại, và dữ liệu
ngoài miền Wikipedia. Mọi kết luận chỉ có giá trị trên tập validation của UIT-ViQuAD~2.0,
với một seed huấn luyện và cấu hình được mô tả.

\subsection{Đóng góp chính}
\label{sec:contrib}

\begin{enumerate}
  \item \textbf{Đưa PhoBERT vào đọc hiểu trích xuất mà không cần Java.} PhoBERT chỉ có
        tokenizer ``chậm'', không trả về vị trí ký tự (offset). Nhóm cài một lớp bọc tự dựng
        offset từ bước tách từ, được kiểm thử bằng bất biến ``giải mã nhãn phải ra lại đúng
        đáp án vàng'' (Mục~\ref{sec:segtok}).
  \item \textbf{Một phép đo chẩn đoán không cần chạy mô hình}: với mỗi câu hỏi, đo xem biên
        đáp án vàng có trùng biên từ của bộ tách từ không, và tính \emph{trần EM} mà một mô
        hình word-level hoàn hảo đạt được. Phép đo chạy trên toàn bộ \NValAns{} câu có đáp án
        của validation, không phải trên vài chục câu (Mục~\ref{sec:alignment}).
  \item \textbf{So sánh có kiểm soát} PhoBERT và XLM-R trên cùng dữ liệu, cùng lịch học, cùng
        quy trình chọn epoch; chênh lệch được kiểm định bằng McNemar và bootstrap trên từng
        câu thay vì so hai con số tổng (Mục~\ref{sec:paired}).
  \item \textbf{Một bộ stress-test kiểm định được}, dựng từ validation thay vì tự viết câu
        hỏi: mỗi mục hoặc là câu gốc do người gán, hoặc là câu bị biến đổi bằng luật mà nhãn
        mới suy ra được từ nhãn cũ, kèm chính câu gốc để so theo cặp. Bất biến ``mọi mục đều
        chấm được'' được kiểm bằng máy trước khi chấm mô hình (\StressN{} mục, \StressViol{}
        vi phạm; Mục~\ref{sec:audit}).
  \item \textbf{Tái lập được}: mọi con số trong báo cáo được sinh từ tệp
        \code{results/*.json} kèm commit; chính các bảng và con số trong văn bản của báo cáo
        này được sinh tự động từ các tệp đó (Phụ lục~\ref{app:trace}).
\end{enumerate}

\subsection{Cấu trúc báo cáo}

Mục~\ref{sec:theory} trình bày cơ sở lý thuyết. Mục~\ref{sec:method} mô tả phương pháp:
pipeline, cách đưa PhoBERT vào bài toán trích xuất, các phép đo chẩn đoán và giao thức thực
nghiệm. Mục~\ref{sec:experiments} trình bày kết quả và trả lời lần lượt bốn câu hỏi nghiên
cứu. Mục~\ref{sec:conclusion} kết luận và nêu giới hạn. Phụ lục cung cấp hướng dẫn tái lập
và bảng truy vết số liệu.
